% Źródła: Eves s. 304--311 (§7.6, §7.7).

\section{Punkty idealne i ultraidealne} % lang-pl
\section{Punti ideali e ultraideali} % lang-it
\label{sekcja_punkty_idealne}

Na płaszczyźnie euklidesowej dwie proste albo się przecinają, albo są równoległe; na hiperbolicznej dochodzi trzeci przypadek -- proste hiperrównoległe. % lang-pl
W tej sekcji zobaczymy, jak Hilbert opisze ten trzeci przypadek i jak uzupełnić płaszczyznę fikcyjnymi punktami, tak jak w geometrii euklidesowej dokłada się punkty w nieskończoności. % lang-pl
Nel piano euclideo due rette o si intersecano o sono parallele; in quello iperbolico si aggiunge un terzo caso -- le rette iperparallele. % lang-it
In questa sezione vedremo come Hilbert descriverà questo terzo caso e come completare il piano con punti fittizi, come nella geometria euclidea si aggiungono i punti all'infinito. % lang-it

\begin{theorem}
[Hilberta] % lang-pl
[di Hilbert] % lang-it
	Dwie proste hiperrównoległe mają dokładnie jedną wspólną prostopadłą. % lang-pl
	Due rette iperparallele hanno esattamente una perpendicolare comune. % lang-it
\end{theorem}

Dowód Hilberta jest konstruktywny: Eves \cite[s. 305--306]{eves1_1972}. % lang-pl
Eves zwraca uwagę, że konstrukcje w geometrii Łobaczewskiego (równoległych przez dany punkt, wspólnej równoległej do dwóch przecinających się prostych) bynajmniej nie są oczywiste \cite[s. 306]{eves1_1972}. % lang-pl
La dimostrazione di Hilbert è costruttiva: Eves \cite[p. 305--306]{eves1_1972}. % lang-it
Eves osserva che le costruzioni nella geometria di Lobachevskij (delle parallele per un punto dato, della parallela comune a due rette incidenti) non sono affatto evidenti \cite[p. 306]{eves1_1972}. % lang-it
\index[persons]{Hilbert, David}%

\begin{definition}
[punkt idealny, punkt ultraidealny] % lang-pl
[punto ideale, punto ultraideale] % lang-it
\index{punkt!idealny} % lang-pl
\index{punkt!ultraidealny} % lang-pl
\index{punto!ideale} % lang-it
\index{punto!ultraideale} % lang-it
	Każdej rodzinie półprostych równoległych przypisujemy wspólny punkt idealny. % lang-pl
	Każdej rodzinie prostych hiperrównoległych, prostopadłych do jednej prostej, przypisujemy wspólny punkt ultraidealny. % lang-pl
	A ogni famiglia di semirette parallele assegniamo un punto ideale comune. % lang-it
	A ogni famiglia di rette iperparallele, perpendicolari a una stessa retta, assegniamo un punto ultraideale comune. % lang-it
\end{definition}

Eves \cite[s. 306]{eves1_1972}. % lang-pl
Eves \cite[p. 306]{eves1_1972}. % lang-it

Wszystko to da się zobaczyć na jednym rysunku: zwykłe punkty płaszczyzny Łobaczewskiego to punkty wewnątrz ustalonego koła $K$, punkty idealne leżą na okręgu $K$, ultraidealne na zewnątrz (łącznie z nieskończonymi), a proste to proste płaszczyzny rozszerzonej przecinające wnętrze koła; prosta stowarzyszona z punktem ultraidealnym to jego biegunowa względem $K$ \cite[s. 306--307]{eves1_1972}. % lang-pl
Arthur Cayley nazwie okrąg $K$ \emph{absolutem} płaszczyzny (zobacz sekcję \ref{sekcja_beltrami_cayley_klein}). % lang-pl
Eves widzi w tym na razie tylko wygodny rysunek, który pozwala zapamiętać wiele faktów; dopiero po przygotowaniu tła posłuży on do pokazania związku z geometrią rzutową \cite[s. 307]{eves1_1972}. % lang-pl
Tutto ciò si può vedere in un unico disegno: i punti ordinari del piano di Lobachevskij sono i punti interni a un cerchio fissato $K$, quelli ideali stanno sulla circonferenza $K$, quelli ultraideali all'esterno (compresi quelli all'infinito), e le rette sono le rette del piano esteso che attraversano l'interno del cerchio; la retta associata a un punto ultraideale è la sua polare rispetto a $K$ \cite[p. 306--307]{eves1_1972}. % lang-it
Arthur Cayley chiamerà la circonferenza $K$ l'\emph{assoluto} del piano (si veda la sezione \ref{sekcja_beltrami_cayley_klein}). % lang-it
Eves vi vede per il momento solo un comodo disegno che permette di ricordare molti fatti; solo dopo aver preparato lo sfondo servirà a mostrare il legame con la geometria proiettiva \cite[p. 307]{eves1_1972}. % lang-it
\index{absolut} % lang-pl
\index{assoluto} % lang-it

W tym obrazie każde dwie proste przecinają się w punkcie zwykłym, idealnym lub ultraidealnym, ale nie każde dwa punkty wyznaczają prostą \cite[s. 307]{eves1_1972}. % lang-pl
Nel disegno ogni coppia di rette si interseca in un punto ordinario, ideale o ultraideale, ma non ogni coppia di punti determina una retta \cite[p. 307]{eves1_1972}. % lang-it

\begin{theorem}
	Symetralne boków dowolnego trójkąta przecinają się w jednym punkcie -- zwykłym, idealnym albo ultraidealnym. % lang-pl
	Gli assi dei lati di un triangolo qualsiasi si incontrano in un unico punto -- ordinario, ideale o ultraideale. % lang-it
\end{theorem}

Dowód (rozbity na trzy przypadki): Eves \cite[s. 308--311]{eves1_1972}. % lang-pl
Dimostrazione (divisa in tre casi): Eves \cite[p. 308--311]{eves1_1972}. % lang-it

Zadania: sprawdzenie trzynastu faktów w modelu absolutu -- Eves \cite[s. 307--308]{eves1_1972}. % lang-pl
Esercizi: verifica di tredici fatti nel modello dell'assoluto -- Eves \cite[p. 307--308]{eves1_1972}. % lang-it
